\section{Related Work}
\label{sec:related_work}

\subsection{Spatial Understanding with VLMs}
Spatial understanding underpins robotic manipulation by recovering geometric relationships. Approaches overcome VLM limitations through two main strategies.
Tool-augmented methods~\cite{ding2024open6dor, qin2025robofactory, zhudexflywheel, tan2025roboos} delegate VLMs to orchestrate external visual foundation models or geometric engines. SpaceTools~\cite{chen2025spacetools} combines dual-interaction reinforcement learning with geometric reasoning, while TIGER~\cite{han2025tiger} synthesizes code-based geometric routines. These externalize computation but add overhead and require module coupling.
3D-data-driven fine-tuning~\cite{cheng2024spatialrgpt, song2025robospatial, daxberger2025mm, raysat, wang2025spatial457} strengthens spatial representations through pseudo-3D or real-world 3D data. RoboRefer~\cite{zhouroborefer} integrates depth encoders, SpatialBot~\cite{cai2025spatialbot} applies progressive spatial curricula, and SpatialVLM~\cite{chen2024spatialvlm} leverages large-scale 3D spatial QA datasets. Fine-tuning incurs substantial data and training costs.
RoboStream contrasts by anchoring visual evidence directly to explicit 3D geometry via STF-Tokens, requiring no fine-tuning and maintaining persistent spatial representations across sequential steps without external modules or retraining.

\subsection{VLMs for Robotic Manipulation}
Robotic manipulation bridges semantic reasoning with execution along three complementary axes.
Logic-grounded task planning~\cite{huang2022language, singh2023progprompt, zengsocratic, huang2023grounded} decomposes instructions into structured sequences. SayCan~\cite{brohan2023saycan} scores action feasibility via affordance models, while Code as Policies~\cite{liang2023code} synthesizes executable programs for long-horizon scheduling. These methods prioritize logical structure but treat execution steps independently.
Geometry-aware constraint generation~\cite{yu2023language, sharma2022correcting, yuan2025robopoint, huang2025rekep, huang2025a3vlm, fang2024moka} grounds planning in spatial constraints. VoxPoser~\cite{huang2023voxposer} synthesizes 3D value maps for zero-shot planning, CoPa~\cite{huang2024copa} defines constraints over geometric keypoints, and SoFar~\cite{qisofar} maps language to 6-DoF poses. However, these approaches reconstruct geometric constraints at each step without persistent tracking.
Closed-loop reasoning~\cite{driess2023palm, du2023vision, zhou2025code, liuinteractive} enhances adaptability through perceptual feedback. Inner Monologue~\cite{huang2023inner} implements self-correcting loops, and VILA~\cite{hu2024look} reasons over image sequences for physical intuition. Yet these methods remain reactive, recovering from failures without reasoning about action-induced state changes or preconditions.
RoboStream unifies geometry and causality through a CSTG that records action-state transitions, enabling VLMs to maintain causal consistency while reasoning about occluded states and external disturbances.

\subsection{Long-Horizon Task Planning}
Long-horizon manipulation poses a distinct challenge, requiring systems to translate high-level instructions into coherent action sequences while maintaining consistency across extended temporal spans. Dynamic visual conditions and partial observability amplify this challenge regarding causal consistency across sequential decisions~\cite{hanrobocerebra}.
Hierarchical decomposition~\cite{shi2024yell, belkhale2024rt, lihamster} organizes multi-step execution by encoding intermediate representations. HiRobot~\cite{shi2025hi} segments instructions into atomic commands, DexVLA~\cite{wen2025dexvla} couples VLM planning with diffusion-based policies, and RoboMatrix~\cite{mao2024robomatrix} implements three-tier hierarchical control. LoHoVLA~\cite{yang2025lohovla} aligns semantic intent via natural-language sub-goals. These strategies organize execution without tracking how past actions reshape the environment.
Chain-of-thought methods~\cite{wei2022chain, mu2023embodiedgpt, zawalski2025robotic, zhang2025learning} strengthen planning via explicit future-state reasoning. CoT-VLA~\cite{zhao2025cot} reasons over latent visual subgoals, ManualVLA~\cite{gu2025manualvla} generates executable visual guidance, and GR-3~\cite{cheang2025gr} learns flow-matching trajectories. While improving short-horizon decisions, these methods lack persistent memory to track action consequences.
RoboStream addresses this gap through a CSTG recording actions, environmental consequences, and external disturbances, enabling VLMs to trace state evolution and reason causally about how past actions constrain future possibilities under partial observability.